\section{引言}
\label{sec:introduction}

低空无人机在巡检、监测、配送和应急响应等任务中具有部署灵活和快速到达的优势。在建筑密集、禁飞区交错且风险分布不均的低空空间中，路径规划需要同时权衡障碍物避让、禁飞区规避、高度边界、转弯平滑性、风险暴露和能耗，而非仅寻找一条短路径。已有综述将 UAV 路径规划概括为图搜索、采样、数学规划、智能优化、学习和混合等多类路线，并指出复杂环境中的三维可行性保持与多约束折中仍是核心问题 \citep{Kumar2025UAVReview,Ghambari2024UAVSurvey}。

围绕复杂约束下的搜索效率与路径质量，同行已从初始化、候选生成、更新和局部引导等环节提出改进。Darlan 等将 A* 引导的交叉、变异和自适应多项式变异嵌入进化算法 \citep{Darlan2026DomainSpecificEA}；Niu 等在连续蚁群优化中结合随机游走和自适应路径点修复 \citep{Niu2025ContinuousACO}；Chen 等面向多 UAV 威胁环境，将偏置采样、候选评价和路径重构整合到 RRT* 框架 \citep{Chen2025BiasedSampling}；Wang 等在多目标水母搜索中利用 RRT 初始化和类别最优个体引导改善初始路径与搜索压力 \citep{Wang2024JellyfishUAV}；Meng 等则通过自适应算子选择和知识融合增强进化多任务路径搜索 \citep{Meng2025EvoMultitasking}。在局部方向引导方面，Khatib 提出了经典人工势场思想 \citep{Khatib1986PotentialField}，Liu 等进一步以旋转势场和参数学习服务多 UAV 协同规划 \citep{Liu2025RotationAPF}。这些研究说明领域知识可以有效参与搜索，但本文所面对的痛点不同：对于单条强约束 B 样条路径，惩罚项和 Deb 规则只能在候选评价后识别不可行性；而若将局部方向强制施加给全种群，又可能干扰基础进化搜索。因此，需要一种能够在评价前提供局部约束方向、按群体状态稀疏介入，并以保守规则限制替换范围的候选生成机制。

针对上述不足，本文提出约束可行性场驱动的自适应进化搜索方法（Constraint Viability Field guided Adaptive Evolution, CVF-AE）。本文不以 CVF 替代基础 AE 或构造新的目标函数，而是将其设计为候选生成中的局部修正机制：约束感知初始化和参考弱引导首先帮助种群进入潜在可行走廊；形成、保持和恢复三种状态根据群体可行性与搜索进展决定何时介入；对少量近可行候选，CVF 将障碍物、禁飞区、风险、高度、边界和转弯信息编码为有界控制点增量；候选仅在 Deb 顺序改善且含惩罚综合适应度不增加时才被保留。由此，初始化、状态调度、场构建和保守接收共同构成一个在不接管全局搜索前提下前移约束信息的闭环。基于这一思路，本文的贡献如下：

\begin{itemize}
    \item 提出面向 B 样条控制点编码的 CVF 表示，将障碍物、禁飞区、高度与边界等局部约束信息，以及风险和转弯等质量信息，映射为控制点层面的有界方向增量，使局部信息在候选评价前参与搜索，而不只通过目标项和惩罚项事后体现。
    \item 设计以“进入--介入--移动”为主线的稀疏候选增强框架：约束感知初始化提供弱先验，三状态调度和近可行触发决定介入时机与对象，基础 AE 方向与 CVF 方向在受限步长内融合，从而保留全局进化搜索的主体地位。
    \item 给出“接受--替换”的保守边界：CVF 候选经过边界投影后，须同时满足局部 Deb 优先和含惩罚综合适应度不增条件才可替代基础候选；稀疏可行性保持仅作用于少量候选，以避免将局部修正扩展为全种群强制修复。
\end{itemize}

本文其余部分组织如下：第 \ref{sec:problem-formulation} 节给出路径编码、目标函数、约束违反量和可行性准则；第 \ref{sec:method} 节介绍 CVF-AE 的候选生成、约束状态识别、CVF 构建和稀疏触发机制；第 \ref{sec:results-discussion} 节依次说明整体实验设置、主对比、消融、显著性和开销分析，并在相应小节内讨论结果含义与方法边界；最后给出结论。
